\loeskopf{Einheit 3 von 5 · Normalform und p-q-Formel (Nr. 17--28)}

\erg{17}{a) nur $x^2$ \quad b) Normalform \quad c) Produkt $= 0$ \quad
d) Klammer im Quadrat \quad e) Normalform}

\erg{18}{a) $p = 10$, $q = 3$ \quad b) $p = -4$, $q = 9$ \quad c) $p = 7$, $q = -5$
\quad d) $p = -9$, $q = -2$ \quad e) $p = -1$, $q = 6$ \quad
f) geordnet $x^2 - 8x + 15 = 0$; $p = -8$, $q = 15$}

\erg{19}{a) nein \quad b) ja, durch $3$ \quad c) nein \quad d) ja, durch $-1$ \quad
e) ja, durch $5$}

\erg{20}{a) $x_1 = -3$, $x_2 = -5$ \quad b) $x_1 = -3$, $x_2 = -7$ \quad
c) $x_1 = -1$, $x_2 = -5$ \quad d) $x_1 = -4$, $x_2 = -8$ \quad
e) $x_1 = 3$, $x_2 = -7$ \quad f) $x_1 = 6$, $x_2 = 4$ \quad
g) $x_1 = 7$, $x_2 = -5$ \quad h) $p$ halbe $= 2{,}5$; $x_1 = -2$, $x_2 = -3$}

\erg{21}{a) unter der Wurzel $0$, genau eine Lösung: $x = 7$ \quad
b) unter der Wurzel $-5$, keine Lösung \quad
c) unter der Wurzel $5$, zwei Lösungen: $x_1 = 4 + \sqrt{5} \approx 6{,}24$,
$x_2 = 4 - \sqrt{5} \approx 1{,}76$}

\erg{22}{a) $x^2 + 5x - 24 = 0$; $x_1 = 3$, $x_2 = -8$ \quad
b) $x^2 - 7x - 18 = 0$; $x_1 = 9$, $x_2 = -2$ \quad
c) $x^2 + 6x + 9 = 4x + 21$, also $x^2 + 2x - 12 = 0$;
$x_{1,2} = -1 \pm \sqrt{13}$, also $x_1 \approx 2{,}61$, $x_2 \approx -4{,}61$}

\erg{23}{a) $x^2 + 7x + 6 = 0$; $x_1 = -1$, $x_2 = -6$ \quad
b) $x^2 - 6x + 8 = 0$; $x_1 = 4$, $x_2 = 2$ \quad
c) mal $(-1)$: $x^2 - 2x - 8 = 0$; $x_1 = 4$, $x_2 = -2$}

\erg{24}{a) $x_1 = 5$, $x_2 = 4$; Probe mit $x = 5$: $25 - 45 + 20 = 0$ (wA) \quad
b) $x_1 = 2$, $x_2 = -5$; Probe mit $x = -5$: $25 - 15 - 10 = 0$ (wA)}

\erg{25}{a) durch $2$ teilen: $x^2 + 7x + 12 = 0$; $x_1 = -3$, $x_2 = -4$ \quad
b) ordnen und mal $(-1)$: $x^2 + 4x - 12 = 0$; $x_1 = 2$, $x_2 = -6$ \quad
c) $x_{1,2} = 5 \pm \sqrt{2}$, also $x_1 \approx 6{,}41$, $x_2 \approx 3{,}59$}

\erg{26}{Fehler: das Vorzeichen von $-\frac{p}{2}$. Wegen $p = -8$ ist
$-\frac{p}{2} = +4$, nicht $-4$. Richtig: $x_{1,2} = 4 \pm \sqrt{16 - 12} = 4 \pm 2$,
also $x_1 = 6$, $x_2 = 2$. \quad a) $x_1 = 9$, $x_2 = 3$}

\erg{27}{a) Die Formel ist für die Normalform $x^2 + px + q = 0$ hergeleitet. Steht
eine andere Zahl vor $x^2$, sind die abgelesenen $p$ und $q$ nicht die der Formel --
erst teilen, dann ablesen. \quad b) Es gibt keine Lösung: Eine Wurzel aus einer
negativen Zahl gibt es nicht, also liefert die Formel keinen Wert. \quad
c) Nein. Das $\pm$ liefert höchstens zwei Werte -- zwei, wenn unter der Wurzel etwas
Positives steht, einen, wenn dort null steht, sonst keinen.}

\erg{28}{a) Wurzelziehen \quad b) Nullprodukt (ausklammern) \quad c) Lösungsformel
\quad d) Wurzelziehen (rückwärts) \quad e) Nullprodukt}
